\chapter{String Algorithms}\label{ch:strings}

\begin{center}\itshape
When a match fails, you have not learnt nothing:\\
you have learnt what you just read, and the pattern already knows what that means.\\
\upshape(the principle of Knuth, Morris and Pratt, in our words)
\end{center}

In the summer of 1969 James Morris was writing a text editor for the CDC 6400 computer, and he wanted to search for a string without ever backing up in the file he was reading. He found a way. Several months later, as the paper that eventually made it famous tells, other programmers had made ``gratuitous fixes'' to his routine, which they found too complicated to understand,
and turned it into a shambles. Independently, early in 1970, Donald Knuth learnt of a surprising theorem of Stephen Cook: every language recognised by a two-way deterministic pushdown automaton, a strange machine that can move back and forth on its input, can be recognised in linear time by a computer. Knuth worked through Cook's construction for a simple automaton, abstracted the mechanism, and obtained a linear-time pattern matcher. He showed it to Vaughan Pratt, who
made its running time independent of the size of the alphabet; and when Pratt described the result to Morris, Morris recognised it as his own. The three published together in 1977. It is a rare case where an algorithm came out of a theorem of automata theory, from someone who had never seen the code, and the code from someone who had never seen the theorem.

This chapter is about searching in texts, which today means searching in web pages, in program sources, in network packets for viruses, and in genomes of three billion letters. We start with one pattern: the Knuth--Morris--Pratt algorithm, which keeps what it has learnt, the $Z$-algorithm, which compares the text with itself, and the Karp--Rabin algorithm, which uses the fingerprints of Chapter~\ref{ch:randomised}. Then many patterns at once, with the automaton of Aho and Corasick.
Then we index the text instead of the pattern: suffix arrays sort all the suffixes of a text, and the Burrows--Wheeler transform turns them into a compressed index that can be searched without decompression. The chapter ends with a sketch that counts the distinct words of an enormous stream in a few kilobytes, the HyperLogLog algorithm promised in Chapter~\ref{ch:randomised}.

Throughout, strings are indexed from $0$: $w=w[0]w[1]\cdots w[|w|-1]$, and $w[i..j)$ is the factor $w[i]\cdots w[j-1]$. The pattern $P$ has length $m$, the text $T$ has length $n$, and both are written with letters of an alphabet $\Sigma$ of size $\sigma$.

\begin{think}
You read a text letter by letter, looking for the pattern \texttt{abacabab}. You have just matched \texttt{abacaba} and the next letter of the text is \texttt{c}, not \texttt{b}. Without looking back at the text, what is the longest beginning of the pattern that you may still have matched? How would you prepare, in advance, the answers to all questions of this kind?
\end{think}

\section{Borders and the Knuth--Morris--Pratt algorithm}\label{sec:kmp}

\begin{definition}
A \emph{border} of a string $w$ is a string that is both a proper prefix and a proper suffix of $w$ (the empty string is always one). For $1\le i\le m$, the \emph{failure function} of $P$ is $\pi(i)$, the length of the longest border of $P[0..i)$.
\end{definition}

A border of length $b$ of $w$ is the same thing as a \emph{period} $|w|-b$: $w[j]=w[j+|w|-b]$ whenever both sides are defined. In the notation of Tome~C1, Chapter~4, the correlation set $O(X,X)$ of a word $X$, which decides Penney's game, is the set of the lengths of its borders, together with $|X|$; computing it is exactly what the next proposition does.

\begin{lemma}\label{lem:borderchain}
\begin{enumerate}
\item The borders of a nonempty string $w$ are its longest border $u$ and the borders of $u$. Hence the lengths of the borders of $P[0..i)$ are $\pi(i)>\pi(\pi(i))>\cdots>0$.
\item For a letter $a$, the nonempty borders of $wa$ are the strings $va$ where $v$ is a border of $w$ with $w[|v|]=a$.
\end{enumerate}
\end{lemma}
\begin{proof}
(1) A border of $u$ is a prefix of a prefix and a suffix of a suffix of $w$, and shorter than $w$. Conversely let $v$ be a border of $w$ with $|v|<|u|$. Both $v$ and $u$ are prefixes of $w$, so $v$ is a prefix of $u$; both are suffixes of $w$, so $v$ is a suffix of $u$; and $v\ne u$. So $v$ is a border of $u$. Iterating gives the chain.

(2) Let $v$ be a string with $|v|<|w|$. Then $va$ is a prefix of $wa$ if and only if $v$ is a prefix of $w$ and $w[|v|]=a$; and $va$ is a suffix of $wa$ if and only if $v$ is a suffix of $w$. Finally $va$ is proper in $wa$ if and only if $v$ is proper in $w$.
\end{proof}

So $\pi(i+1)$ is $1+|v|$ for the longest border $v$ of $P[0..i)$ followed in $P$ by the letter $P[i]$, if there is one, and $0$ otherwise; and the borders of $P[0..i)$ are visited in decreasing order by iterating $\pi$.

\begin{proposition}\label{prop:failure}
The following computes $\pi(1),\dots,\pi(m)$ with fewer than $2m$ comparisons of letters: set $\pi(1)=0$ and $k=0$; for $i=1,\dots,m-1$, while $k>0$ and $P[k]\ne P[i]$ replace $k$ by $\pi(k)$; then if $P[k]=P[i]$ increase $k$ by $1$; set $\pi(i+1)=k$.
\end{proposition}
\begin{proof}
At the start of step $i$, $k=\pi(i)$ (true for $i=1$). By Lemma~\ref{lem:borderchain}, the loop tries the borders of $P[0..i)$ in decreasing order of length, $k,\pi(k),\dots$, and stops at the first one, of length $k$, with $P[k]=P[i]$, or at $k=0$; then $\pi(i+1)=k+1$ if $P[k]=P[i]$ and $0$ otherwise, which is what is stored. For the count: every unsuccessful comparison inside the loop decreases $k$ (as $\pi(k)<k$), each step ends with one more comparison, and $k$ increases by at most $1$ per step and never goes below $0$.
So there are at most $m-1$ decreases, and at most $(m-1)+(m-1)$ comparisons.
\end{proof}

\begin{theorem}[Knuth, Morris, Pratt]\label{thm:kmp}
Set $q=0$. For $j=0,\dots,n-1$: while $q=m$, or $q>0$ and $P[q]\ne T[j]$, replace $q$ by $\pi(q)$; then if $P[q]=T[j]$ increase $q$ by $1$; if now $q=m$, report an occurrence of $P$ ending at position $j$. This reports all occurrences of $P$ in $T$, with fewer than $2n$ comparisons of letters; together with Proposition~\ref{prop:failure}, the time is $O(n+m)$.
\end{theorem}
\begin{proof}
Let $s_j$ be the length of the longest prefix of $P$ that is a suffix of $T[0..j]$, and $s_{-1}=0$. We show that after step $j$, $q=s_j$; then occurrences, which are the $j$ with $s_j=m$, are exactly reported.

The prefixes of $P$ that are suffixes of $T[0..j)$ are $P[0..s_{j-1})$ and its borders: a shorter such prefix is a prefix and a suffix of $P[0..s_{j-1})$, by the argument of Lemma~\ref{lem:borderchain}(1). A prefix $ua$ of $P$, with $a$ a letter, is a suffix of $T[0..j]$ if and only if $u$ is a prefix of $P$ of length less than $m$ that is a suffix of $T[0..j)$ and $P[|u|]=a=T[j]$.
So $s_j$ is $1+|u|$ for the longest $u$ in the list $P[0..s_{j-1})$, its longest border, and so on, with $|u|<m$ and $P[|u|]=T[j]$, and $0$ if none exists. This is what the loop computes, starting from $q=s_{j-1}$. The count of comparisons is as in Proposition~\ref{prop:failure}: $q$ increases at most $n$ times in all, and each unsuccessful comparison decreases it.
\end{proof}

The algorithm never moves backwards in the text, as Morris wanted, and the time it spends on one letter of the text is not bounded by a constant, but its average is. A variant that precomputes all transitions (Exercise~\ref{ex:kmpautomaton}) spends exactly one step per letter, at the cost of $O(m\sigma)$ preprocessing. In practice, for a single pattern in natural text, the Boyer--Moore algorithm of the same year, which compares from the right end of the pattern and can skip ahead by up to $m$ letters, is often faster still, since it typically looks at only a fraction of the text.

\section{The $Z$-algorithm}\label{sec:zalgo}

A second linear-time method compares a string with itself. For a string $S$ of length $N$ and $1\le i<N$, let $Z[i]$ be the length of the longest common prefix of $S$ and $S[i..N)$. Each $i$ with $Z[i]>0$ defines a \emph{$Z$-box} $[i,i+Z[i])$, a window that repeats the beginning of $S$.

\begin{theorem}[the $Z$-algorithm]\label{thm:zalgo}
The array $Z[1],\dots,Z[N-1]$ can be computed with $O(N)$ comparisons, as follows. Keep the $Z$-box $[l,r)$ with the largest $r$ found so far (initially $l=r=0$). For $i=1,\dots,N-1$: if $i<r$ and $Z[i-l]<r-i$, set $Z[i]=Z[i-l]$; otherwise, starting from $z=\max(0,r-i)$, increase $z$ while $i+z<N$ and $S[z]=S[i+z]$, set $Z[i]=z$, and if $i+z>r$ replace $[l,r)$ by $[i,i+z)$.
\end{theorem}
\begin{proof}
Since $S[l..r)=S[0..r-l)$, for $l<i<r$ we have $S[i..r)=S[i-l..r-l)$. If $Z[i-l]<r-i$, the common prefix of $S$ with $S[i-l..N)$ ends before position $r-l$ with a mismatch $S[Z[i-l]]\ne S[i-l+Z[i-l]]$, and the right-hand letter equals $S[i+Z[i-l]]$ because $i+Z[i-l]<r$; and the letters before agree. So $Z[i]=Z[i-l]$.
Otherwise $Z[i]\ge r-i$ (the copied part matches up to $r$, when $i<r$), so starting the explicit comparisons at $z=\max(0,r-i)$ is correct. Each successful explicit comparison with $i+z\ge r$ moves $r$ forward, and $r\le N$; each $i$ has at most one unsuccessful comparison. So there are at most $2N$ comparisons.
\end{proof}

To find $P$ in $T$, apply this to $S=P\#T$, where $\#$ is a letter that occurs in neither: $P$ occurs at position $i$ of $T$ if and only if $Z[m+1+i]=m$. The $Z$-algorithm was made popular by Dan Gusfield's textbook of 1997; the same idea of reusing a box of known matches drives Glenn Manacher's linear-time palindrome algorithm of 1975.

\section{Fingerprints: Karp and Rabin}\label{sec:karprabin}

Richard Karp and Michael Rabin proposed in 1987 a method that replaces comparisons of strings by comparisons of numbers. Encode the letters as elements of $\mathbb F_p$ for a prime $p\ge\sigma$, choose $x\in\mathbb F_p$ uniformly at random, and define the \emph{fingerprint} of a string $w$ of length $m$ as
\[
H(w)=\sum_{i=0}^{m-1}w[i]\,x^{m-1-i}\in\mathbb F_p.
\]
The fingerprints $H_i=H(T[i..i+m))$ of all the windows satisfy $H_{i+1}=\bigl(H_i-T[i]x^{m-1}\bigr)x+T[i+m]$, so all of them cost $O(n)$ operations in $\mathbb F_p$ once $x^{m-1}$ is known.

\begin{theorem}[Karp--Rabin]\label{thm:karprabin}
The Karp--Rabin algorithm reports every $i$ with $H_i=H(P)$. It reports every occurrence of $P$, and the expected number of reported windows that are not occurrences is at most $(n-m+1)(m-1)/p$. If each reported window is verified letter by letter, the algorithm becomes always correct, and its expected running time is $O(n+m+m\cdot\mathrm{occ}+nm^2/p)$, where $\mathrm{occ}$ is the number of occurrences; with $p\ge nm$ this is $O(n+m(\mathrm{occ}+1))$.
\end{theorem}
\begin{proof}
If $T[i..i+m)=P$ the fingerprints are equal. If not, $H_i-H(P)$ is the value at $x$ of a nonzero polynomial of degree at most $m-1$, which vanishes with probability at most $(m-1)/p$ (Proposition~\ref{prop:fingerprint}). Linearity of expectation over the $n-m+1$ windows gives the bound, and each false report costs $O(m)$ to verify.
\end{proof}

The great strength of fingerprints is their flexibility: they compare windows of different texts, find repeated factors by hashing all windows of a given length, and extend to two dimensions, to search for a block of pixels in an image. Their weakness is the randomness: without verification, a false match occurs with small but positive probability, and an adversary who knows $x$ can construct one.

\section{Many patterns: Aho and Corasick}\label{sec:aho}

In 1975 Alfred Aho and Margaret Corasick, at Bell Laboratories, published \emph{Efficient string matching: an aid to bibliographic search}: they wanted to find, in one pass over a text, all occurrences of any of a large set of keywords. Their algorithm is Knuth--Morris--Pratt for a set of patterns. It became the engine of the Unix command \texttt{fgrep} and, since then, of virus scanners and network intrusion detectors.

Let $D$ be a finite set of nonempty patterns of total length $M$. The \emph{trie} of $D$ is the tree whose nodes are the prefixes of the patterns, the root being the empty string, with an edge labelled $a$ from $u$ to $ua$. We identify each node with its string. For a node $v\ne\varepsilon$, the \emph{failure link} $f(v)$ is the longest proper suffix of $v$ that is a node, and the \emph{output link} $o(v)$ is the longest proper suffix of $v$ that is a pattern, if there is one (Figure~\ref{fig:aho}).

\begin{figure}[ht]
\centering
\fig{ch11-aho}
\caption{The Aho--Corasick automaton of $\{\texttt{he},\texttt{she},\texttt{his},\texttt{hers}\}$, the example of the 1975 paper. Black edges form the trie; double circles mark the patterns; dashed arrows are the failure links that do not go to the root $0$. Reading \texttt{she}, the automaton reaches node $5$, whose failure link leads to node $2$, so both \texttt{she} and \texttt{he} are reported.}
\label{fig:aho}
\end{figure}

\begin{theorem}[Aho--Corasick]\label{thm:aho}
Assume that the child of a node along a given letter can be found in constant time (for instance with hash tables). The failure and output links can be computed in $O(M)$ time, and then all occurrences of all patterns of $D$ in a text $T$ of length $n$ are reported in $O(n+\mathrm{occ})$ time, where $\mathrm{occ}$ is the number of occurrences.
\end{theorem}
\begin{proof}
\emph{Search.} Keep a node $v$; for each letter $a=T[j]$, while $v\ne\varepsilon$ and $v$ has no child along $a$, replace $v$ by $f(v)$; then if $v$ has a child along $a$, move to it. By the argument of Theorem~\ref{thm:kmp}, after step $j$ the node $v$ is the longest suffix of $T[0..j]$ that is a node: the suffixes of a node $v$ that are nodes are $v,f(v),f(f(v)),\dots$ (as in Lemma~\ref{lem:borderchain}(1)), and $ua$ is a node and a suffix of $T[0..j]$ if and only if $u$ is a node, a suffix of $T[0..j)$, and has a child along $a$.
A pattern that is a suffix of $T[0..j]$ is a node, hence a suffix of $v$; so the patterns ending at $j$ are $v$ itself if it is a pattern, then $o(v)$, $o(o(v))$, and so on, each found in constant time. The depth of $v$ increases by at most $1$ per letter and decreases at each failure step, so there are at most $n$ failure steps.

\emph{Preprocessing.} Process the nodes in order of increasing depth. The children of the root have $f=\varepsilon$. For a node $u=va$ with $v\ne\varepsilon$, the proper suffixes of $u$ that are nodes are $\varepsilon$ and the strings $wa$ with $w$ a proper suffix of $v$ that is a node with a child along $a$ (a prefix of a node is a node). So $f(u)$ is found by walking $w=f(v),f(f(v)),\dots$ until some $w$ has a child $wa$, giving $f(u)=wa$, or until $w=\varepsilon$ without such a child, giving $f(u)=\varepsilon$.
Then $o(u)=f(u)$ if $f(u)$ is a pattern and $o(u)=o(f(u))$ otherwise. For the cost, follow one pattern $P$ along its path $u_1,u_2,\dots,u_{|P|}$ in the trie: the walk computing $f(u_{i+1})$ starts at depth $|f(u_i)|$, each step decreases the depth, and $|f(u_{i+1})|\le|f(u_i)|+1$. So the steps along the path are at most $|P|$, and assigning each node to one pattern whose path contains it, the total is at most $M$.
\end{proof}

The output term $\mathrm{occ}$ cannot be avoided: there can be many more occurrences than letters of text (Exercise~\ref{ex:ahooutputs}). Precomputing the full transition function, as in Exercise~\ref{ex:kmpautomaton}, gives a deterministic automaton with one step per letter of text: this is the form used in hardware for network scanning.

\section{Suffix arrays}\label{sec:suffixarray}

The algorithms so far preprocess the pattern. When the text is fixed and the patterns arrive one by one, as for a genome that is searched with millions of short reads, one should preprocess the text instead. Peter Weiner showed in 1973 that the \emph{suffix tree} of a text, the trie of all its suffixes with chains of unary nodes compressed, can be built in linear time; but suffix trees typically use ten to twenty bytes per letter.
In 1990 Udi Manber and Gene Myers proposed a leaner structure, which keeps only the order.

Let $S$ be a string of length $N$ whose last letter is a sentinel $\$$, smaller than all other letters and occurring nowhere else. Its \emph{suffix array} $\mathrm{SA}$ lists the starting positions of the suffixes $S[i..N)$ in increasing lexicographic order. The \emph{LCP array} gives, for $1\le k<N$, the length $\mathrm{LCP}[k]$ of the longest common prefix of the suffixes at ranks $k-1$ and $k$. For \texttt{banana\$}, $\mathrm{SA}=(6,5,3,1,0,4,2)$, listing \texttt{\$}, \texttt{a\$}, \texttt{ana\$}, \texttt{anana\$}, \texttt{banana\$}, \texttt{na\$}, \texttt{nana\$}.

\begin{lemma}[sandwich]\label{lem:sandwich}
If $x\le y\le z$ in lexicographic order, then $\mathrm{lcp}(x,z)\le\mathrm{lcp}(y,z)$ and $\mathrm{lcp}(x,z)\le\mathrm{lcp}(x,y)$.
\end{lemma}
\begin{proof}
Let $\ell=\mathrm{lcp}(x,z)$ and $c$ the common prefix of length $\ell$. If $y$ did not begin with $c$, let $t<\ell$ be the first position where $y$ differs from $c$ (or $y$ is a proper prefix of $c$). If $y$ is shorter, or $y[t]<c[t]$, then $y<x$; if $y[t]>c[t]$, then $y>z$. Both contradict the hypothesis, so $y$ begins with $c$.
\end{proof}

\begin{theorem}[prefix doubling]\label{thm:doubling}
The suffix array of $S$ can be computed in $O(N\log N)$ time.
\end{theorem}
\begin{proof}
For $k\ge0$ let $\rho_k(i)$ be the rank of the factor $S[i..\min(i+2^k,N))$ among all the distinct factors of this form, and set $\rho_k(i)=-1$ for $i\ge N$. The ranks $\rho_0$ are those of the letters, found by sorting. We claim that $S[i..i+2^{k+1})$ (truncated at $N$) is smaller than $S[j..j+2^{k+1})$ if and only if the pair $(\rho_k(i),\rho_k(i+2^k))$ is smaller than $(\rho_k(j),\rho_k(j+2^k))$ in lexicographic order of pairs, and that they are equal if and only if the pairs are.
Indeed, compare the first halves: if they differ, their order decides. If they are equal and one of them is truncated, it contains the sentinel, which occurs once, so $i=j$. Otherwise both first halves have length $2^k$, and the order is decided by the second halves, the empty string (rank $-1$) being smallest.

So $\rho_{k+1}$ is obtained from $\rho_k$ by sorting $N$ pairs of integers in $[-1,N)$, which a two-pass radix sort does in $O(N)$ time (Chapter~\ref{ch:sorting}). Once $2^k\ge N$, the factors are whole suffixes, which are all distinct, so $\rho_k$ is a permutation and $\mathrm{SA}$ is its inverse. This happens after $\lceil\log_2N\rceil$ rounds.
\end{proof}

Linear-time constructions exist; the first three were found independently in 2003, by Kärkkäinen and Sanders, by Ko and Aluru, and by Kim, Sim, Park and Park, and the induced-sorting algorithm of Nong, Zhang and Chan (2009) is the one used in practice. With the suffix array, a pattern $P$ is found by binary search: the suffixes that begin with $P$ form an interval of ranks, and two binary searches, each comparing $P$ with $O(\log N)$ suffixes, find its ends in $O(m\log N)$ time. Many questions about repeats need the LCP array too.

\begin{theorem}[Kasai, Lee, Arimura, Arikawa, Park]\label{thm:kasai}
Given $S$ and $\mathrm{SA}$, the LCP array can be computed in $O(N)$ time.
\end{theorem}
\begin{proof}
Let $\mathrm{rk}$ be the inverse permutation of $\mathrm{SA}$, and for each position $i$ let $h(i)$ be the length of the common prefix of $S[i..N)$ and the suffix just before it in the array (and $h(i)=0$ if $\mathrm{rk}(i)=0$); then $\mathrm{LCP}[\mathrm{rk}(i)]=h(i)$. We claim that $h(i+1)\ge h(i)-1$. Suppose $h(i)=\ell\ge1$, and let $j$ be the position just before $i$ in the array, so $S[j..N)<S[i..N)$ with a common prefix of length $\ell$.
Since $\$$ occurs once and $\ell\ge1$, neither suffix is $\$$, so removing their common first letter gives $S[j+1..N)<S[i+1..N)$ with a common prefix of length $\ell-1$. The suffix just before $S[i+1..N)$ in the array lies between these two, so by Lemma~\ref{lem:sandwich} its common prefix with $S[i+1..N)$ has length at least $\ell-1$.

So compute $h(0),h(1),\dots,h(N-1)$ in this order, comparing letters from position $\max(h(i-1)-1,0)$ onwards instead of from $0$. Each successful comparison increases the current value, which never exceeds $N$ and decreases by at most $1$ per position; each position has one unsuccessful comparison. So there are $O(N)$ comparisons.
\end{proof}

The longest repeated factor of $S$ has length $\max_k\mathrm{LCP}[k]$; the number of distinct factors is $\sum$ of the suffix lengths minus $\sum_k\mathrm{LCP}[k]$ (Exercise~\ref{ex:distinctsub}); and the longest common factor of two strings comes from one suffix array of both (Exercise~\ref{ex:lcsubstring}).

\section{The Burrows--Wheeler transform}\label{sec:bwt}

In 1983, while working at AT\&T Bell Laboratories, David Wheeler discovered a strange permutation of a text. Write all the rotations of the text in a list, sort the list, and keep only the last column: the result is a permutation of the text, which can be undone, and which tends to put equal letters together. Wheeler did not publish it. Eleven years later, on 10 May 1994, he and Michael Burrows described it in a research report of the Digital Systems Research Center, \emph{A block-sorting lossless data compression algorithm}. Within a few years it was the heart of the compressor \texttt{bzip2}; within fifteen it was the index of the programs that align DNA reads to the human genome.

Let $S$ end with the sentinel $\$$ as before, and sort its $N$ rotations; since $\$$ occurs once, sorting the rotations is the same as sorting the suffixes, and the rotation at rank $k$ starts at $\mathrm{SA}[k]$. Let $F$ and $L$ be the first and last columns of the sorted list (Figure~\ref{fig:bwt}). The string $L$, with $L[k]=S[\mathrm{SA}[k]-1]$ (and $L[k]=\$$ if $\mathrm{SA}[k]=0$), is the \emph{Burrows--Wheeler transform} of $S$; $F$ is just the sorted list of the letters.
For a letter $c$, let $C[c]$ be the number of letters of $S$ smaller than $c$, and $\mathrm{rank}_c(k)$ the number of occurrences of $c$ in $L[0..k)$.

\begin{figure}[ht]
\centering
\fig{ch11-bwt}
\caption{The sorted rotations of \texttt{banana\$}. The first column $F$ (blue) is sorted; the last column $L$ (red), read downwards, is the transform \texttt{annb\$aa}. Each row is a suffix of the text followed by the beginning of the text, and the suffix array is the second column.}
\label{fig:bwt}
\end{figure}

\begin{lemma}[last to first]\label{lem:lf}
For every letter $c$, the occurrences of $c$ in $L$, read from top to bottom, are the same letters of $S$ as the occurrences of $c$ in $F$, in the same order. Consequently, if $L[k]=c$, the rotation that starts one position earlier than the rotation of row $k$ is in row $\mathrm{LF}(k)=C[c]+\mathrm{rank}_c(k)$.
\end{lemma}
\begin{proof}
Let $k<k'$ be rows with $L[k]=L[k']=c$, with rotations $X=Yc$ and $X'=Y'c$. The rotations one step earlier are $cY$ and $cY'$, which begin with $c$. Two distinct rotations of $S$ already differ among their first $N-1$ letters (otherwise they would have the same multiset of letters outside the last one, hence the same last letter, and be equal; and the rotations are distinct because $\$$ occurs once). So $X<X'$ implies $Y<Y'$, hence $cY<cY'$:
the order of the rows with $L=c$ is the order of the rows with $F=c$ that contain the same letters. The rows beginning with $c$ form the block of ranks $C[c],\dots,C[c]+(\text{number of }c)-1$, and within it the row of $cY$ is the $\mathrm{rank}_c(k)$-th, counting from $0$.
\end{proof}

\begin{theorem}\label{thm:bwtinvert}
\begin{enumerate}
\item $S$ can be recovered from $L$ in $O(N)$ time, assuming a constant alphabet.
\item Given $L$, $C$ and a structure answering the queries $\mathrm{rank}_c(k)$ in constant time, the number of occurrences of a pattern $P$ of length $m$ in $S$ is computed with $O(m)$ rank queries (\emph{backward search}).
\end{enumerate}
\end{theorem}
\begin{proof}
(1) Compute $C$ by counting, and all values $\mathrm{rank}_{L[k]}(k)$ in one scan of $L$. Row $0$ is the rotation $\$S[0..N-1)$, which starts at position $N-1$; its last letter is $S[N-2]$. By Lemma~\ref{lem:lf}, iterating $k\mapsto\mathrm{LF}(k)$ from $k=0$ visits the rotations starting at $N-2,N-3,\dots,0$, and the letters $L[k]$ read along the way are $S[N-2],S[N-3],\dots,S[0]$, the text backwards.

(2) The rows whose rotations begin with a string $Q$ form an interval $[s,e)$, since the rotations are sorted. For a letter $c$, the rows beginning with $cQ$ are the images under $\mathrm{LF}$ of the rows in $[s,e)$ whose last letter is $c$; by Lemma~\ref{lem:lf} these images keep their order and are consecutive, so they form the interval $[C[c]+\mathrm{rank}_c(s),\,C[c]+\mathrm{rank}_c(e))$.
Start with the interval of all rows, for the empty string, and prepend the letters of $P$ from the last to the first; after $m$ steps the interval has as many rows as $P$ has occurrences, and two rank queries were used per step.
\end{proof}

This is the \emph{FM-index} of Paolo Ferragina and Giovanni Manzini (2000). Rank queries can be answered in constant time with a small amount of extra space, by storing counts at regular checkpoints, and the positions of the occurrences are recovered by storing a sample of the suffix array and walking with $\mathrm{LF}$ to the nearest sampled row. Meanwhile $L$ itself compresses well: the rows that begin with \texttt{he~} in an English text end mostly with \texttt{t}, so $L$ has long runs of equal letters, which a move-to-front transform followed by entropy coding exploits (the notebook measures it). The result is an index of the human genome that fits in a few gigabytes of memory, used since 2009 by read aligners such as Bowtie and BWA.

\section{Counting distinct strings}\label{sec:hll}

How many distinct search queries did a search engine receive today, or how many distinct $31$-letter factors does a genome contain? Storing them all to count them is out of the question. In 1985 Philippe Flajolet and Nigel Martin observed that a hash function makes the stream look random, and that the most unusual hash value seen is a good measure of how many distinct values there were.

Let $h$ map the universe $U$ of strings to infinite sequences of independent fair bits, and let $\varrho(x)\ge1$ be the position of the first $1$ in $h(x)$. We assume that $h$ is a truly random function: in practice it is a good hash function, and the analysis is a model of it. The \emph{register} of a stream is $R=\max\varrho(x)$ over the stream; repeated elements do not change it.

\begin{proposition}\label{prop:fmregister}
If the stream contains $F$ distinct elements, then for every $c\ge0$,
\[
\mathbb P\bigl(2^R\ge2^{c+1}F\bigr)\le2^{-c},\qquad\mathbb P\bigl(2^R\le2^{-c}F\bigr)\le\exp(-2^{c-1}).
\]
If $h(x)$ is truncated to $\log_2|U|$ bits, the register needs $\log_2\log_2|U|+O(1)$ bits.
\end{proposition}
\begin{proof}
For each distinct element $x$, $\mathbb P(\varrho(x)\ge r)=2^{1-r}$ for integers $r\ge1$, independently over $x$. For real $r$, the union bound gives $\mathbb P(R\ge r)=\mathbb P(R\ge\lceil r\rceil)\le F2^{1-\lceil r\rceil}\le F2^{1-r}$, and independence gives $\mathbb P(R<r)\le(1-2^{1-\lceil r\rceil})^F\le\exp(-F2^{1-\lceil r\rceil})\le\exp(-F2^{-r})$, since $\lceil r\rceil<r+1$.
With $r=\log_2F+c+1$ the first bound is $2^{-c}$. The event $2^R\le2^{-c}F$ is $R\le\log_2F-c$, which implies $R<\log_2F-c+1$, of probability at most $\exp(-F2^{-\log_2F+c-1})=\exp(-2^{c-1})$. Finally $R\le\log_2|U|+1$, a number of $\log_2\log_2|U|+O(1)$ bits.
\end{proof}

One register estimates $F$ within a constant factor, with constant probability: far too crude, but in a few bits. The \emph{HyperLogLog} algorithm of Flajolet, \'Eric Fusy, Olivier Gandouet and Fr\'ed\'eric Meunier (2007) uses $\mu$ registers: the first bits of $h(x)$ choose a register $j$, and $\varrho$ of the remaining bits updates register $M_j$. The estimate is
\[
E=\alpha_\mu\,\mu^2\Bigl(\sum_{j=1}^\mu2^{-M_j}\Bigr)^{-1},
\]
a harmonic mean of the $2^{M_j}$ scaled by a constant $\alpha_\mu\approx0.7213/(1+1.079/\mu)$ that corrects the bias. They proved, by a delicate analysis with Mellin transforms that we only quote, that the relative standard error is close to $1.04/\sqrt\mu$. With $\mu=2048$ registers of $5$ bits, about $1.3$ kilobytes, distinct counts up to billions are estimated within about $2\%$. Two sketches of two streams combine by taking the maximum register by register, which gives the sketch of their union (Exercise~\ref{ex:hllmerge}); this is why HyperLogLog is built into many databases.

\begin{beauty}
A string algorithm is a way of not reading the same thing twice. Knuth, Morris and Pratt store, in the borders of the pattern, everything a failed comparison can teach; Aho and Corasick do the same for a whole dictionary; the suffix array sorts once all the questions that will ever be asked of a text; and the Burrows--Wheeler transform, a permutation found in 1983 and kept in a drawer for eleven years, is at the same time a compressor and an index,
because the order that brings similar contexts together is also the order that answers searches. HyperLogLog goes furthest: it reads each word once and forgets it, remembering only how lucky the luckiest hash was.
\end{beauty}

\section{Exercises}

\begin{exercise}[$\star$]\label{ex:pifunction}
Compute the failure function of $P=\texttt{abacabab}$ and its periods. Run Theorem~\ref{thm:kmp} on $T=\texttt{abacabacabab}$, and answer the question of the box at the beginning of the chapter.
\end{exercise}

\begin{exercise}[$\star\star$]\label{ex:periodicity}
\emph{Periodicity lemma (weak form).} Show that if $p$ and $q$ are periods of $w$ with $p+q\le|w|$, then $\gcd(p,q)$ is also a period of $w$. (Fine and Wilf proved in 1965 that $p+q-\gcd(p,q)\le|w|$ suffices, and that this is optimal.)
\end{exercise}

\begin{exercise}[$\star$]\label{ex:kmpautomaton}
Define $\delta(q,a)$ for $0\le q\le m$ and $a\in\Sigma$ as the length of the longest prefix of $P$ that is a suffix of $P[0..q)a$. Show that $\delta(q,a)=q+1$ if $q<m$ and $P[q]=a$, that otherwise $\delta(q,a)=\delta(\pi(q),a)$ for $q\ge1$ and $\delta(0,a)=0$, and deduce an $O(m\sigma)$ construction of an automaton that searches $T$ with exactly one transition per letter.
\end{exercise}

\begin{exercise}[$\star$]\label{ex:zperiod}
Show that the smallest period of a string $S$ of length $N$ is the least $p\ge1$ with $p+Z[p]=N$, or $N$ if there is none. Deduce a linear algorithm that decides whether $S$ is a power $u^k$ with $k\ge2$.
\end{exercise}

\begin{exercise}[$\star$]\label{ex:krprime}
How large must $p$ be for the Karp--Rabin algorithm without verification to make no false report at all, with probability at least $1-1/n$? How many bits does $p$ have for a text of $10^9$ letters and a pattern of $1000$ letters, and what can one do if machine words have $64$ bits?
\end{exercise}

\begin{exercise}[$\star$]\label{ex:ahooutputs}
Let $D=\{\texttt{a},\texttt{aa},\dots,\texttt{a}^k\}$ and $T=\texttt{a}^n$ with $n\ge k$. Compute $M$ and the number of occurrences, and show that $\mathrm{occ}$ can be of order $n\sqrt M$.
\end{exercise}

\begin{exercise}[$\star\star$]\label{ex:distinctsub}
Let $S'$ be a string of length $n$ and $S=S'\$$. Show that the number of distinct nonempty factors of $S'$ is $\frac{n(n+1)}2-\sum_{k}\mathrm{LCP}[k]$, and check it for \texttt{banana}.
\end{exercise}

\begin{exercise}[$\star\star$]\label{ex:lcsubstring}
Let $X$ and $Y$ be strings and $S=X\#Y\$$ with $\#,\$$ new letters. Show that the length of the longest common factor of $X$ and $Y$ is the largest $\mathrm{LCP}[k]$ over the ranks $k$ such that one of the suffixes at ranks $k-1$ and $k$ starts in $X$ and the other in $Y$.
\end{exercise}

\begin{exercise}[$\star$]\label{ex:bwtbanana}
Starting only from $L=\texttt{annb\$aa}$, recover \texttt{banana\$} with Lemma~\ref{lem:lf}, and count the occurrences of \texttt{ana} by backward search.
\end{exercise}

\begin{exercise}[$\star$]\label{ex:hllmerge}
Show that the HyperLogLog registers of the concatenation of two streams are the maxima of their registers, register by register. Can the size of the intersection of two streams be estimated from their two sketches?
\end{exercise}

\section{Hints}

\begin{hint}{ex:pifunction}
Use Lemma~\ref{lem:borderchain}. A period of $P$ is $m$ minus the length of a border.
\end{hint}
\begin{hint}{ex:periodicity}
With $p>q$, show that $p-q$ is a period, distinguishing positions $i<q$ and $i\ge q$; then imitate Euclid.
\end{hint}
\begin{hint}{ex:kmpautomaton}
This is Lemma~\ref{lem:borderchain}(2) and the proof of Theorem~\ref{thm:kmp}, with $T[0..j)$ replaced by $P[0..q)$.
\end{hint}
\begin{hint}{ex:zperiod}
$p$ is a period if and only if $S[p..N)=S[0..N-p)$.
\end{hint}
\begin{hint}{ex:krprime}
Union bound over the windows.
\end{hint}
\begin{hint}{ex:ahooutputs}
$\texttt{a}^\ell$ occurs $n-\ell+1$ times.
\end{hint}
\begin{hint}{ex:distinctsub}
Every factor is a prefix of a suffix. Count, for each suffix in the order of the array, the prefixes that are not prefixes of an earlier suffix.
\end{hint}
\begin{hint}{ex:lcsubstring}
A common factor is a common prefix of a suffix starting in $X$ and a suffix starting in $Y$. Use Lemma~\ref{lem:sandwich}.
\end{hint}
\begin{hint}{ex:bwtbanana}
$F=\texttt{\$aaabnn}$, and $C[\$]=0$, $C[\texttt a]=1$, $C[\texttt b]=4$, $C[\texttt n]=5$.
\end{hint}
\begin{hint}{ex:hllmerge}
Each register is a maximum over the elements sent to it.
\end{hint}

\section{Solutions}

\begin{solution}{ex:pifunction}
\idea Extend borders one letter at a time, falling back along the chain of borders when the next letter does not fit.
\writeup For the prefixes $\texttt a,\texttt{ab},\dots,\texttt{abacabab}$ the longest borders have lengths $\pi=0,0,1,0,1,2,3,2$: for instance $\texttt{abacaba}$ has the border \texttt{aba}, and then $P[3]=\texttt c\ne\texttt b=P[7]$, so we fall back to $\pi(3)=1$, the border \texttt a, followed by $P[1]=\texttt b$: $\pi(8)=2$. The borders of $P$ are \texttt{ab} and the empty string, so its periods are $8-2=6$ and $8$.
On $T=\texttt{abacabacabab}$ the matched length $q$ goes $1,2,3,4,5,6,7$ on the first seven letters; the eighth letter \texttt c mismatches $P[7]=\texttt b$, so $q$ falls to $\pi(7)=3$, and $P[3]=\texttt c$ matches: $q=4$. Then \texttt{abab} brings $q$ to $8$ at the last letter: one occurrence, at position $4$. This answers the opening question: after matching \texttt{abacaba} and reading \texttt c, the longest prefix still matched is \texttt{abac}, found without looking back at the text.
\end{solution}

\begin{solution}{ex:periodicity}
\idea Two periods can be subtracted as long as the string is long enough, which is Euclid's algorithm on periods.
\writeup Suppose $p>q$; we show that $p-q$ is a period, that is $w[i]=w[i+p-q]$ for $0\le i<|w|-(p-q)$. If $i\ge q$: $w[i]=w[i-q]$ by period $q$, and $w[i-q]=w[i-q+p]$ by period $p$, as $i-q+p<|w|$. If $i<q$: then $i+p<p+q\le|w|$, so $w[i]=w[i+p]$ by period $p$, and $w[i+p]=w[i+p-q]$ by period $q$.
Now $p-q$ and $q$ are periods with $(p-q)+q\le|w|$, and their gcd is $\gcd(p,q)$. By induction on $p+q$ (the case $p=q$ being trivial), $\gcd(p,q)$ is a period.
\end{solution}

\begin{solution}{ex:kmpautomaton}
\idea The automaton is the search algorithm with all its fall-backs precomputed.
\writeup A string $ua$ is a prefix of $P$ that is a suffix of $P[0..q)a$ if and only if $u$ is a prefix of $P$ of length less than $m$, a suffix of $P[0..q)$, and $P[|u|]=a$ (Lemma~\ref{lem:borderchain}(2), in the same form as in the proof of Theorem~\ref{thm:kmp}). If $q<m$ and $P[q]=a$, the choice $u=P[0..q)$ is allowed and is the longest, so $\delta(q,a)=q+1$.
Otherwise $u$ must be a prefix of $P$ that is a proper suffix of $P[0..q)$, that is a border of $P[0..q)$. For $q\ge1$ these are $P[0..\pi(q))$ and its borders (Lemma~\ref{lem:borderchain}(1)), which are exactly the prefixes of $P$ that are suffixes of $P[0..\pi(q))$; and these are the candidates in the definition of $\delta(\pi(q),a)$. So $\delta(q,a)=\delta(\pi(q),a)$. For $q=0$ the only candidate is the empty $u$, which fails since $P[0]\ne a$: $\delta(0,a)=0$.
As $\pi(q)<q$, the table can be filled in increasing order of $q$ with $O(1)$ work per entry, $O(m\sigma)$ in all. Then $q_j=\delta(q_{j-1},T[j])$ is the number $s_j$ of the proof of Theorem~\ref{thm:kmp}, computed with one lookup per letter of text, and an occurrence ends at $j$ when $q_j=m$.
\end{solution}

\begin{solution}{ex:zperiod}
\idea A period is a self-overlap, and $Z$ measures self-overlaps.
\writeup By definition, $p$ is a period of $S$ if and only if $S[i]=S[i+p]$ for $0\le i<N-p$, that is $S[p..N)=S[0..N-p)$, that is the common prefix of $S$ and $S[p..N)$ is all of $S[p..N)$: $Z[p]=N-p$. So the smallest period $p_0$ is the least $p$ with $p+Z[p]=N$, or $N$, which is always a period.
We claim that $S$ is a power $u^k$ with $k\ge2$ if and only if $p_0<N$ and $p_0\mid N$. If $p_0<N$ divides $N$, then $S=(S[0..p_0))^{N/p_0}$. Conversely if $S=u^k$ with $k\ge2$, then $|u|$ is a period, $|u|\le N/2$ and $p_0\le|u|$, so $p_0+|u|\le N$; by Exercise~\ref{ex:periodicity}, $\gcd(p_0,|u|)$ is a period, hence equals $p_0$ by minimality, so $p_0\mid|u|\mid N$ and $p_0<N$. Computing $Z$ and then $p_0$ takes linear time.
\end{solution}

\begin{solution}{ex:krprime}
\idea The bound of Theorem~\ref{thm:karprabin} is the expected number of false reports, hence also a bound on the probability that there is one.
\writeup By the union bound over windows, the probability of at least one false report is at most $(n-m+1)(m-1)/p<nm/p$. It is at most $1/n$ if $p\ge n^2m$. For $n=10^9$ and $m=1000$ this is $p\ge10^{21}$, a prime of $70$ bits. With $64$-bit words one can use two independent fingerprints, with two random points modulo two primes near $2^{61}$ (or modulo the same prime): a false report needs both to fail, which has probability at most $(m/p)^2$ per window, about $10^{-31}$ here, and the union bound over $10^9$ windows gives about $10^{-22}$.
\end{solution}

\begin{solution}{ex:ahooutputs}
\idea Nested patterns overlap everywhere, so occurrences outnumber letters.
\writeup $M=1+2+\dots+k=\frac{k(k+1)}2$. The pattern $\texttt a^\ell$ occurs at the $n-\ell+1$ positions $0,\dots,n-\ell$ of $T$, so $\mathrm{occ}=\sum_{\ell=1}^k(n-\ell+1)=kn-\frac{k(k-1)}2\ge\frac{kn}2$. As $k^2\ge\frac{k(k+1)}2=M$, we have $k\ge\sqrt M$ and $\mathrm{occ}\ge\frac12n\sqrt M$: no algorithm that lists the occurrences can run in time $O(n+M)$, and the term $\mathrm{occ}$ of Theorem~\ref{thm:aho} is needed.
\end{solution}

\begin{solution}{ex:distinctsub}
\idea Each suffix, taken in sorted order, contributes its prefixes minus those already seen, and those are exactly the ones it shares with its predecessor.
\writeup Every nonempty factor of $S'$ is a nonempty prefix, without $\$$, of a suffix of $S$. The suffix of rank $0$ is $\$$, which has no such prefix. For $k\ge1$, the suffix of rank $k$ is $S'[i..n)\$$ with $i=\mathrm{SA}[k]$, and it has $n-i$ nonempty prefixes without $\$$. By Lemma~\ref{lem:sandwich}, such a prefix is a prefix of a suffix of smaller rank if and only if it is a prefix of the suffix of rank $k-1$, that is if and only if its length is at most $\mathrm{LCP}[k]$ (a common prefix of two distinct suffixes cannot contain $\$$, which occurs once).
So the suffix of rank $k$ contributes $n-i-\mathrm{LCP}[k]$ new factors, and summing over $k\ge1$ gives $\sum_{i=0}^{n-1}(n-i)-\sum_k\mathrm{LCP}[k]=\frac{n(n+1)}2-\sum_k\mathrm{LCP}[k]$, where $\mathrm{LCP}[1]=0$ since the suffix of rank $0$ is $\$$. For \texttt{banana}, the suffixes \texttt{a}, \texttt{ana}, \texttt{anana}, \texttt{banana}, \texttt{na}, \texttt{nana} (followed by $\$$) have LCP values $0,1,3,0,0,2$ with their predecessors, of sum $6$, so there are $21-6=15$ distinct factors.
\end{solution}

\begin{solution}{ex:lcsubstring}
\idea A longest common factor is shared by two suffixes from different sides; between them in the array there is an adjacent pair from different sides that shares at least as much.
\writeup Let $\ell$ be the largest LCP over adjacent mixed pairs. Such a common prefix cannot contain $\#$ or $\$$, which occur once each, so it lies inside $X$ and inside $Y$: it is a common factor, and the longest common factor has length at least $\ell$. Conversely, a common factor $W$ is a common prefix of a suffix $x$ starting in $X$ and a suffix $y$ starting in $Y$, say with $x<y$. Walking through the array from $x$ to $y$, some adjacent pair $(u,v)$ consists of one suffix from each side, with $x\le u<v\le y$.
By Lemma~\ref{lem:sandwich} applied twice, $\mathrm{lcp}(u,v)\ge\mathrm{lcp}(x,v)\ge\mathrm{lcp}(x,y)\ge|W|$. So $\ell\ge|W|$.
\end{solution}

\begin{solution}{ex:bwtbanana}
\idea Walk backwards with the last-to-first map; search backwards with intervals.
\writeup $L=\texttt{annb\$aa}$ gives $F=\texttt{\$aaabnn}$ and $C[\$]=0$, $C[\texttt a]=1$, $C[\texttt b]=4$, $C[\texttt n]=5$. The ranks $\mathrm{rank}_{L[k]}(k)$ for $k=0,\dots,6$ are $0,0,1,0,0,1,2$, so $\mathrm{LF}=(1,5,6,4,0,2,3)$. From row $0$: $L[0]=\texttt a$, go to row $1$: $L[1]=\texttt n$, row $5$: $\texttt a$, row $2$: $\texttt n$, row $6$: $\texttt a$, row $3$: $\texttt b$, row $4$: $\$$. The letters read, \texttt a, \texttt n, \texttt a, \texttt n, \texttt a, \texttt b, spell \texttt{banana} backwards, and the row of $\$$ ends the walk.
Backward search for \texttt{ana}: start with $[0,7)$. Letter \texttt a: $[1+\mathrm{rank}_{\texttt a}(0),\,1+\mathrm{rank}_{\texttt a}(7))=[1,4)$. Letter \texttt n: $[5+\mathrm{rank}_{\texttt n}(1),\,5+\mathrm{rank}_{\texttt n}(4))=[5,7)$. Letter \texttt a: $[1+\mathrm{rank}_{\texttt a}(5),\,1+\mathrm{rank}_{\texttt a}(7))=[2,4)$. Two rows, so \texttt{ana} occurs twice, at positions $\mathrm{SA}[2]=3$ and $\mathrm{SA}[3]=1$.
\end{solution}

\begin{solution}{ex:hllmerge}
\idea Registers are maxima, and a maximum over a union is the maximum of the maxima.
\writeup The register $M_j$ of a stream is the maximum of $\varrho$ over the distinct elements of the stream that the hash sends to register $j$ (and $0$ if there are none). The elements of the concatenated stream sent to $j$ are those of the first stream and those of the second, so its register is the maximum of the two registers. So sketches of different machines or days merge exactly, and the merged sketch estimates the size of the union. For the intersection there is no such operation; one can use inclusion--exclusion, $|A\cap B|=|A|+|B|-|A\cup B|$, but the error of each estimate is a fraction of $|A|$ or $|A\cup B|$, so small intersections of large sets are estimated very poorly.
\end{solution}

\section*{Sources and further reading}

\begin{itemize}
\item \emph{The stories.} D.~E.~Knuth, J.~H.~Morris and V.~R.~Pratt, \emph{Fast pattern matching in strings}, SIAM Journal on Computing 6 (1977), 323--350, whose Section~7, \emph{Historical remarks}, tells the story of the opening; S.~A.~Cook, \emph{Linear time simulation of deterministic two-way pushdown automata}, Information Processing 71 (1972), 75--80. M.~Burrows and D.~J.~Wheeler, \emph{A block-sorting lossless data compression algorithm}, Research Report 124, Digital Systems Research Center, 10 May 1994, whose introduction dates Wheeler's discovery to 1983 at AT\&T Bell Laboratories.
\item R.~S.~Boyer and J.~S.~Moore, \emph{A fast string searching algorithm}, Communications of the ACM 20 (1977), 762--772; R.~M.~Karp and M.~O.~Rabin, \emph{Efficient randomized pattern-matching algorithms}, IBM Journal of Research and Development 31 (1987), 249--260; A.~V.~Aho and M.~J.~Corasick, \emph{Efficient string matching: an aid to bibliographic search}, Communications of the ACM 18 (1975), 333--340; N.~J.~Fine and H.~S.~Wilf, \emph{Uniqueness theorems for periodic functions}, Proceedings of the American Mathematical Society 16 (1965), 109--114; G.~Manacher, \emph{A new linear-time ``on-line'' algorithm for finding the smallest initial palindrome of a string}, Journal of the ACM 22 (1975), 346--351.
\item P.~Weiner, \emph{Linear pattern matching algorithms}, 14th Annual Symposium on Switching and Automata Theory (1973), 1--11; U.~Manber and G.~Myers, \emph{Suffix arrays: a new method for on-line string searches}, SIAM Journal on Computing 22 (1993), 935--948 (conference version 1990); T.~Kasai, G.~Lee, H.~Arimura, S.~Arikawa and K.~Park, \emph{Linear-time longest-common-prefix computation in suffix arrays and its applications}, Combinatorial Pattern Matching 2001, Lecture Notes in Computer Science 2089, 181--192; G.~Nong, S.~Zhang and W.~H.~Chan, \emph{Linear suffix array construction by almost pure induced-sorting}, Data Compression Conference 2009, 193--202.
\item P.~Ferragina and G.~Manzini, \emph{Opportunistic data structures with applications}, 41st IEEE Symposium on Foundations of Computer Science (2000), 390--398; B.~Langmead, C.~Trapnell, M.~Pop and S.~L.~Salzberg, \emph{Ultrafast and memory-efficient alignment of short DNA sequences to the human genome}, Genome Biology 10 (2009), R25; H.~Li and R.~Durbin, \emph{Fast and accurate short read alignment with Burrows--Wheeler transform}, Bioinformatics 25 (2009), 1754--1760.
\item P.~Flajolet and G.~N.~Martin, \emph{Probabilistic counting algorithms for data base applications}, Journal of Computer and System Sciences 31 (1985), 182--209; P.~Flajolet, \'E.~Fusy, O.~Gandouet and F.~Meunier, \emph{HyperLogLog: the analysis of a near-optimal cardinality estimation algorithm}, Analysis of Algorithms 2007, Discrete Mathematics and Theoretical Computer Science Proceedings AH, 127--146.
\item D.~Gusfield, \emph{Algorithms on Strings, Trees, and Sequences}, Cambridge University Press, 1997, the classic reference, with the $Z$-algorithm and the biological motivation; M.~Crochemore, C.~Hancart and T.~Lecroq, \emph{Algorithms on Strings}, Cambridge University Press, 2007; G.~Navarro, \emph{Compact Data Structures}, Cambridge University Press, 2016, for the FM-index and rank structures.
\end{itemize}
